\uuid{Zcyu}
\titre{Système linéaire/Système homogène}
\auteur{hanani ; mzari}
\organisation{exo7}
\licence{CC BY-NC-SA 4.0}
\datecreate{2020-01-20}
\niveau{L1}
\module{Algèbre}
\chapitre{Déterminant, système linéaire}
\sousChapitre{Système linéaire, rang}
\difficulte{3}
\theme{Systèmes d'équations linéaires, Système linéaire, Système homogène}
\exo7qcmid{lille-2.systemes.020}

\contenu{
\begin{question}
\qclassification{200.03}
\qtags{motcle=Système linéaire/Système homogène}
Soit $(\mathtt{S})$ un système à $3$ équations linéaires et $3$ inconnues et $(\mathtt{S}_H)$ le système homogène associé. Quelles sont les assertions vraies ?
\begin{answers}
\bad{$(\mathtt{S}_H)$ admet une infinité de solutions.}
\bad{$(\mathtt{S})$ admet une unique solution.}
\good{Si $X_1$ et $X_2$ sont des solutions de $(\mathtt{S})$, alors $X_1-X_2$ est une solution de $(\mathtt{S}_H)$.}
\bad{$(\mathtt{S})$ admet une infinité de solutions si et seulement si les équations de $(\mathtt{S})$ sont celles de $3$ plans confondus.}
\end{answers}
\begin{explanations}
$(\mathtt{S}_H)$ admet au moins le zéro de l'espace comme solution, mais n'admet pas nécessairement une infinité de solutions. Les équations de $(\mathtt{S})$ étant celles de $3$ plans, $4$ cas sont possibles :
\begin{enumerate}
\item[-] Les $3$ plans sont confondus, dans ce cas, $(\mathtt{S})$ admet une infinité de solutions.
\item[-] Les $3$ plans se coupent en une droite, dans ce cas, $(\mathtt{S})$ admet une infinité de solutions.
\item[-] Les $3$ plans se coupent en un point, dans ce cas, $(\mathtt{S})$ admet une unique solution.
\item[-] L'intersection des  $3$ plans est vide, dans ce cas, $(\mathtt{S})$ n'admet pas de solution.
\end{enumerate}
\end{explanations}
\end{question}
}
